% Abhakseite „Das kann ich“ – Zone nur im Gesamt (\mitzone)
%% TODO Ich-kann-Sätze: Platzhalter wie in den Aufgabendateien
\begin{abhakseite}
\ifdefined\mitzone
\abhakgruppe{Kennst du schon}
\abhak{1}{Achsenspiegelung und Symmetrie ebener Figuren (Spiegelpunkt gleich weit hinter der Achse, Deckabbildung)}
\abhak{2}{Verbindungsvektor bilden und an einem Punkt abtragen (Ortsvektor plus Vektor)}
\abhak{3}{Normalenvektor einer Ebene ablesen und die Lotgerade aufstellen}
\abhak{4}{Schnittpunkt von Gerade und Ebene berechnen (Einsetzen, Parameter bestimmen)}
\abhak{5}{Mittelpunkt einer Strecke berechnen}
\abhak{6}{Beträge von Vektoren vergleichen (gleich lange Richtungsvektoren erkennen oder herstellen)}
\abhak{7}{Verbindungsvektor bilden und an einem Punkt abtragen (Ortsvektor plus Vektor) – Fehler finden}
\abhak{8}{Verbindungsvektor bilden und an einem Punkt abtragen (Ortsvektor plus Vektor) – selbst rechnen}
\fi
\abhakgruppe{Einheit 1 · Punkte spiegeln}
\abhak{9}{Was ist gegeben, was gesucht?}
\abhak{10}{Welche Koordinate wechselt?}
\abhak{11}{Punkte spiegeln}
\abhak{12}{Punkte spiegeln – Fehler finden, Begründen}
\abhakgruppe{Einheit 2 · Spiegelebene und Spiegelgerade bestimmen}
\abhak{13}{Spiegelebene und Spiegelgerade}
\abhak{14}{Spiegelebene und Spiegelgerade – Fehler finden, Begründen}
\abhakgruppe{Einheit 3 · Symmetrieebenen von Körpern}
\abhak{15}{Symmetrieebenen von Körpern}
\abhak{16}{Symmetrieebenen von Körpern – Fehler finden, Begründen}
\end{abhakseite}
